% ==================================================================
% id:       seq-022
% topic:    סדרות
% subtopic: חישוב גבולות בעזרת אריתמטיקה של גבולות
% solution: full
% verified: no
% ==================================================================
\textbf{חשבו את הגבולות הבאים:}
\begin{enumerate}\bMs{2}
\item $\ds\limn\frac{2n^2-3n}{5n^2+1}$
\begin{solbox}
\unverified
נחלק מונה ומכנה ב-$n^2$:
$$\frac{2n^2-3n}{5n^2+1}=\frac{2-\frac3n}{5+\frac{1}{n^2}}\longrightarrow\frac{2-0}{5+0}=\frac25 .$$
הערה: בחישוב הגבול השתמשנו באריתמטיקה של גבולות: הגבול של סכום, של מכפלה ושל מנה של סדרות מתכנסות שווה לסכום, למכפלה ולמנה של הגבולות, כל עוד הגבול של המכנה שונה מ-$0$.
\end{solbox}
\item $\ds\limn\frac{n^3+1}{2n^4-n}$
\begin{solbox}
\unverified
נחלק מונה ומכנה ב-$n^4$:
$$\frac{n^3+1}{2n^4-n}=\frac{\frac1n+\frac{1}{n^4}}{2-\frac{1}{n^3}}\longrightarrow\frac{0+0}{2-0}=0 .$$
הערה: בחישוב הגבול השתמשנו באריתמטיקה של גבולות: הגבול של סכום, של מכפלה ושל מנה של סדרות מתכנסות שווה לסכום, למכפלה ולמנה של הגבולות, כל עוד הגבול של המכנה שונה מ-$0$.
\end{solbox}
\item $\ds\limn\frac{(2n+1)(3n-2)}{n^2+4}$
\begin{solbox}
\unverified
נחלק מונה ומכנה ב-$n^2$:
$$\frac{(2n+1)(3n-2)}{n^2+4}=\frac{\lr{2+\frac1n}\lr{3-\frac2n}}{1+\frac{4}{n^2}}\longrightarrow\frac{2\cdot3}{1}=6 .$$
הערה: בחישוב הגבול השתמשנו באריתמטיקה של גבולות: הגבול של סכום, של מכפלה ושל מנה של סדרות מתכנסות שווה לסכום, למכפלה ולמנה של הגבולות, כל עוד הגבול של המכנה שונה מ-$0$.
\end{solbox}
\item $\ds\limn\lr{\sqrt{n^2+3n}-n}$
\begin{solbox}
\unverified
נכפול ונחלק בצמוד:
$$\sqrt{n^2+3n}-n=\frac{3n}{\sqrt{n^2+3n}+n}=\frac{3}{\sqrt{1+\frac3n}+1}\longrightarrow\frac{3}{1+1}=\frac32 ,$$
כי
$\sqrt{1+\frac3n}=\lr{1+\frac3n}^{1/2}\to1$
(גבול של חזקה).
הערה: בחישוב הגבול השתמשנו באריתמטיקה של גבולות: הגבול של סכום, של מכפלה ושל מנה של סדרות מתכנסות שווה לסכום, למכפלה ולמנה של הגבולות, כל עוד הגבול של המכנה שונה מ-$0$.
\end{solbox}
\item $\ds\limn\lr{\sqrt{n+1}-\sqrt n}$
\begin{solbox}
\unverified
נכפול ונחלק בצמוד:
$$\sqrt{n+1}-\sqrt n=\frac{1}{\sqrt{n+1}+\sqrt n}.$$
מתקיים
$0<\frac{1}{\sqrt{n+1}+\sqrt n}<\frac{1}{\sqrt n}$,
ו-$\frac{1}{\sqrt n}=\lr{\frac1n}^{1/2}\to0$.
לפי כלל הסנדוויץ' הגבול הוא $0$.
\end{solbox}
\item $\ds\limn\lr{\frac{n+1}{2n}}^{\frac{n+1}{n}}$
\begin{solbox}
\unverified
הבסיס
$\frac{n+1}{2n}=\frac12+\frac{1}{2n}$
חיובי ושואף ל-$\frac12>0$,
והמעריך
$\frac{n+1}{n}=1+\frac1n$
שואף ל-$1$.
לפי משפט הגבול של חזקה, הגבול הוא
$\lr{\frac12}^1=\frac12$.
\end{solbox}
\eMs
\end{enumerate}
